% \InsertMark{ztex-4th}{}
\clearpage
\subsection*{\titlepkg{cmd} 库}
\ztex{} 的 \pkg{cmd} 库提供了一种新的命令定义方式: 类似 Python 中的 
\texttt{def fun (\meta{arg-spec}) \marg{code}}; 该库目前很不成熟, 请谨慎使用.


\begin{function}[added=2025-06-19, updated=2025-09-30]{\znewcmd, \zsetcmd, \zgsetcmd}
  \begin{syntax}
    \cs{znewcmd}\meta{cmd}\marg{arg-spec}\marg{code}
  \end{syntax}
  用户可以使用这三个命令创建控制序列, 新创建的控制序列名称为 \meta{cmd}. \meta{arg-spec} 的格式为: \texttt{\meta{var}:\meta{type}=\meta{default}};
  其中 \meta{var} 为\textbf{局部}变量的名称, 可以使用数字, 下划线(但此时需使用 \cs{zcmdvar} 命令进行引用); \meta{type} 用于
  指定变量 \meta{var} 的类型, 可以省略; 目前 \meta{type} 的可选值有 ``\texttt{tl, str, int, fp, clist, dim, \oarg{type}}'',  
  其中 ``\texttt{tl}'' 为默认类型, \texttt{\oarg{type}} 用于表示数组, 数组中元素的类型均为(元素类型必须相同) \meta{type}; \meta{default} 
  用于指定变量 \meta{var} 的默认值, 可以省略; \meta{code} 即为函数体. \par
  \textbf{备注}: 在函数体中, 所有的局部变量均为完全可展的.
  % 数组变量中的每个元素均被包含于 \cmd{\exp_not:n} 中, 普通变量则保留为用户的原始输入形式.
\end{function}



\begin{function}[added=2025-09-30]{\ztex_cmd_create:nnnnn}
  \begin{syntax}
    \cmd{\ztex_cmd_create:nnnnn}
    ~~~~\Arg{cmd}\Arg{arg-parser}\Arg{code}
    ~~~~\Arg{setup}\Arg{boolean}
  \end{syntax}
  此命令用于创建类似 \cmd{\znewcmd}, \cmd{\zsetcmd} 的命令. \meta{cmd} 为函数名, 不包含 ``\texttt{\char92}'';
  \meta{arg-spec} 用于解析命令参数; \meta{code} 为函数的具体执行内容, 可以在其中使用已声明的命令参数; \meta{setup}
  表示 \meta{cmd} 的创建方式, 可以使用 ``\texttt{cs_new:Npn}'' , ``\texttt{cs_set:Npn}'' 等; \meta{boolean}
  为一个布尔值, 用于控制命令参数是否需要计算. \textbf{备注}: 此命令的使用方法可以参考下述 \cmd{\znewcmd} 的定义:
\end{function}
\begin{DocExample}
\cs_set_protected:Npn \znewcmd #1#2#3
  {
    \cs_if_exist:NT {#1}
      {
        \ztex_msg_set:nn {znewcmd@exist}
          {
            command~\string#1~already~exsits!
          }
        \ztex_msg_error:n {znewcmd@exist}
      }
    \exp_args:Ne \ztex_cmd_create:nnnnn {\cs_to_str:N #1}{#2}
      {
        #3
      }{cs_new:Npn}{\c_false_bool}
  }
\end{DocExample}



\begin{function}[added=2025-06-19, EXP]{\fpuse, \intuse, \dimuse, \clistuse}
  \begin{syntax}
    \cs{fpuse}\marg{var}
    \cs{intuse}\marg{var}
    \cs{dimuse}\marg{var}
    \cs{clistuse}\meta{var}\marg{index}
  \end{syntax}
  在 \cs{znewcmd}, \cs{zsetcmd}, \cs{zgsetcmd} 所定义控制序列对应的 \meta{code} 中, 部分的变量并不能直接使用,
  需要使用 \cs{fpuse}, \cs{dimuse} 等命令进行引用.
\end{function}


\begin{function}[added=2025-12-02, EXP]{\cmdvar}
  \begin{syntax}
    \cs{cmdvar}\marg{var}
  \end{syntax}
  此命令用于引用已经声明的变量, 如果被引用的变量含有数字, 横线, 下划线等特殊字符,
  建议使用该命令.
\end{function}

\newgeometry{left=3cm, right=3cm}
最后，我们在这里给出命令 \cmd{\znewcmd} 的一些使用案例, 以供用户参考:
\begin{DocExample}*
\ExplSyntaxOn
\cs_set_eq:NN \tlEQNnTF \tl_if_eq:NnTF
\ExplSyntaxOff
% new command
\znewcmd\CMDA{argA=argA-val, argB:str=argB-val, argC}
  { 
    \tlEQNnTF \argA {argA-val}{argA~EQUALS}{argA~not~EQUALS}\par
    \tlEQNnTF \argB {argB-val}{argB~EQUALS}{argB~not~EQUALS}\par
    \string\argC=\argC\par
  }
\CMDA{argB=argB-val-new}
% set command
\dotfill\par
\zsetcmd\CMDB{
  argA = {``Group variable range Test''},
  argF:fp = 3.1415926,
  argG:int = 100,
  argH:dim = 12pt+1em,
  argI:clist = {AA, BB, CC},
}{
  \string\argF=\fpuse{\argF};
  \string\argG=\intuse\argG;
  \string\argH=\dimuse\argH;
  \string\argI=\clistuse\argI{2}.
  \par\dotfill\par
  Argument of \string\CMDA(local variable test):
  \string\argA=\argA\par
}
\CMDB{argF=6.2830178, argG=200}
% group test
\dotfill\par
\begingroup
\zsetcmd\CMDA{arg-1=aaa}{CODE=\cmdvar{arg-1}}
INNER: \CMDA{};
\endgroup
OUTER: \CMDA{}
% vector type
\dotfill\par
\znewcmd\CMDD{argA:[int]={1, 2, 3, 4}, argB:[str], argC:[tl]}
  { 
    CODE 1=(\argA{1}), (\argA{4})\par
    CODE 2=(\argB{1}), (\argB{-1})\par
    CODE 3=(\argC{1})
  }
\CMDD{argA={5.55, 6, 7, 8}, argB={AAA, BBB, CCC}}
\end{DocExample}
\restoregeometry
